% ============================================================================
%  Viewpoint 5 : O(d)  Cartan 
%  Viewpoint 5 Correction: O(d) Gauge-Fixing is Linearizable via Cartan
%  SCX  — 
%  SCX Gauge Theory — Chinese+English
% ============================================================================
\documentclass[11pt,a4paper]{article}
\usepackage[UTF8]{ctex}
\usepackage{amsmath,amssymb,amsthm}
\usepackage{mathtools}
\usepackage{bm}
\usepackage{booktabs}
\usepackage[margin=2.5cm]{geometry}
\usepackage{xcolor}
\usepackage{hyperref}
\usepackage{enumitem}

\theoremstyle{plain}
\newtheorem{theorem}{Theorem}
\newtheorem{proposition}[theorem]{Proposition}
\newtheorem{lemma}[theorem]{Lemma}

\theoremstyle{definition}
\newtheorem{definition}[theorem]{Definition}

\theoremstyle{remark}
\newtheorem{remark}[theorem]{Remark}
\newtheorem{correction}[theorem]{}

\newcommand{\R}{\mathbb{R}}
\newcommand{\SO}{\mathrm{SO}}
\newcommand{\so}{\mathfrak{so}}
\newcommand{\gauge}{g}
\newcommand{\gaugegroup}{\mathcal{G}}
\newcommand{\Skew}{\mathrm{Skew}}

\title{\textbf{SCX 5}\\[0.3cm]
\large O(d) Cartan vs.\ \\[0.2cm]
\large Viewpoint 5 Correction: Cartan Linearization vs.\ Full Non-Abelian Gauge-Fixing}
\author{SCX  \\ SCX Review Committee}
\date{202672 \\ July 2, 2026}

\begin{document}
\maketitle

% ===================================================================
\section{ () \\ Original Claim (Under Review)}
% ===================================================================

\begin{quote}
\textbf{5 ():}

ACE  O(3)  O(d) 
 O(d) \textbf{}
(non-abelian gauge fixing)  Riemannian Gauss-Newton

\end{quote}

\begin{quote}
\textbf{Viewpoint 5 (Original):}

ACE potentials possess O(3) rotational symmetry, making the gauge group
O(d) rather than the translation group. Because O(d) is a non-abelian
Lie group, gauge fixing requires \textbf{non-abelian gauge fixing}
techniques, implemented as Riemannian Gauss-Newton optimization
on the manifold. This procedure has quadratic convergence.
\end{quote}

% ===================================================================
\section{ \\ Review Verdict}
% ===================================================================

\begin{center}
\fbox{\parbox{0.9\textwidth}{
\centering
\textbf{:  (OVERSTATED)}

\\
(1)  O(d) \\
(2)  Cartan  $\so(d)$ \\
(3) ``''  Yang-Mills \\
\quad Gribov ——\textbf{}\\
(4) Gauss-Newton \textbf{}
}}
\end{center}

\begin{center}
\fbox{\parbox{0.9\textwidth}{
\centering
\textbf{Verdict: OVERSTATED}

The original claim is partially correct in its technical content, but
its formulation is significantly overstated:\\
(1) It conflates the group-theoretic non-abelian property of O(d) with
the need for full non-abelian gauge-fixing machinery;\\
(2) It overlooks the critical fact that Cartan decomposition linearizes
the problem to the Lie algebra $\so(d)$;\\
(3) ``Non-abelian gauge fixing'' in physics specifically denotes Yang-Mills
techniques — ghost fields, Gribov copies, etc. — \textbf{none of which apply};\\
(4) Gauss-Newton has only \textbf{linear convergence} (not quadratic) when
residuals are non-zero.
}}
\end{center}

% ===================================================================
\section{ \\ Why the Original Claim Is Overstated}
% ===================================================================

\subsection{O(d)  ≠ }
\subsection*{O(d) Being Non-Abelian $\neq$ Requiring Non-Abelian Gauge-Fixing}

\noindent
\textbf{.}
O(d)  $R_1, R_2 \in \mathrm{O}(d)$
$R_1 R_2 \neq R_2 R_1$\textbf{
}O(d)  $\so(d)$
\textbf{} $\frac{d(d-1)}{2}$ Cartan 


\medskip
\noindent
\textbf{English.}
O(d) is indeed non-abelian as a group: for $R_1, R_2 \in \mathrm{O}(d)$,
generally $R_1 R_2 \neq R_2 R_1$. However, \textbf{the algebraic property of the
group and the algorithmic complexity of solving the gauge-fixing problem
are distinct concerns}. The crucial fact is: O(d) as a Lie group has Lie
algebra $\so(d)$, which is a \textbf{vector space} (of dimension
$\frac{d(d-1)}{2}$), and Cartan decomposition — the inverse of the
exponential map, i.e., the logarithm — provides a local diffeomorphism
from the group manifold to this vector space.

\subsection{``'' }
\subsection*{What ``Non-Abelian Gauge Fixing'' Means in Physics}

\noindent
\textbf{.}
``'' Yang-Mills  Faddeev-Popov 


\begin{enumerate}[label=(\arabic*)]
  \item \textbf{ (Ghost Fields):} 
  \item \textbf{Gribov  (Gribov Copies):} —
        
  \item \textbf{BRST :} 
  \item \textbf{:}  Jacobian Faddeev-Popov 
\end{enumerate}

\textbf{ ACE  O(d) } ACE 


\medskip
\noindent
\textbf{English.}
In quantum field theory, ``non-abelian gauge fixing'' (e.g., the
Faddeev-Popov procedure in Yang-Mills theory) involves:

\begin{enumerate}[label=(\arabic*)]
  \item \textbf{Ghost Fields:} Anti-commuting scalar fields introduced to
        cancel unphysical gauge degrees of freedom;
  \item \textbf{Gribov Copies:} Non-uniqueness of gauge-fixing conditions —
        the same physical configuration may correspond to multiple gauge-fixed solutions;
  \item \textbf{BRST Symmetry:} A super-symmetry structure relating ghosts
        and gauge fields;
  \item \textbf{Path-Integral Measure:} A Jacobian determinant (Faddeev-Popov
        determinant) is required to correct the integration measure.
\end{enumerate}

\textbf{None of these issues apply to the O(d) gauge-fixing problem for ACE
potentials.} The ACE gauge-fixing is a finite-dimensional parameter estimation
problem, not an infinite-dimensional quantum field theory path-integral problem.

\subsection{Gauss-Newton }
\subsection*{Convergence Order of Gauss-Newton}

\noindent
\textbf{.}
 Gauss-Newton 

\begin{itemize}
  \item Gauss-Newton \textbf{}\textbf{}
        
  \item  ACE ——
        Gauss-Newton \textbf{}
  \item $\|x_{k+1} - x^*\| \leq \rho \|x_k - x^*\|$
         $\rho < 1$  $\rho$  1
\end{itemize}

\medskip
\noindent
\textbf{English.}
The original claim asserts quadratic convergence for Gauss-Newton. This
misreads standard numerical optimization theory:

\begin{itemize}
  \item Gauss-Newton exhibits \textbf{local quadratic convergence} only when
        \textbf{residuals are zero} (zero-residual problem, where the problem
        is exactly satisfied at the solution);
  \item When residuals are non-zero (the typical case for ACE gauge-fixing —
        there is no exact solution that perfectly eliminates all gauge
        redundancy), Gauss-Newton has only \textbf{linear convergence};
  \item The linear convergence rate is governed by the problem's condition
        number: $\|x_{k+1} - x^*\| \leq \rho \|x_k - x^*\|$, where $\rho < 1$
        but $\rho$ may be close to 1.
\end{itemize}

% ===================================================================
\section{ \\ Corrected Claim}
% ===================================================================

\begin{correction}[5 ]
\label{cor:viewpoint5}

ACE  O(3)  O(d) 
 O(d)  \textbf{Cartan }
 $g_m = g_0 \cdot \exp(\xi_m)$
$\xi_m \in \so(d)$  $\so(d)$ 
$\frac{d(d-1)}{2}$ \textbf{}


$\mathcal{O}(d^2)$  Riemannian Gauss-Newton
\textbf{}

\bigskip
\noindent
\textit{ACE potentials possess O(3) rotational symmetry, making the gauge
group O(d) rather than the translation group. While O(d) is non-abelian as
a group, the gauge-fixing problem can be linearized via \textbf{Cartan
decomposition}: writing $g_m = g_0 \cdot \exp(\xi_m)$ with
$\xi_m \in \so(d)$ in the Lie algebra reduces it to a
\textbf{standard linear least-squares problem} on a vector space of
dimension $\frac{d(d-1)}{2}$.

For applications where the log map is valid (near the identity), this
linearization adds only $\mathcal{O}(d^2)$ overhead. The full Riemannian
Gauss-Newton is needed only when far from the solution, and converges
\textbf{linearly} (not quadratically) when residuals are non-zero.}
\end{correction}

% ===================================================================
\section{ \\ Mathematical Justification}
% ===================================================================

\subsection{Cartan }
\subsection*{Cartan Decomposition and the Logarithm Map}

\begin{definition}[Cartan ]
\label{def:cartan}

 $G = \mathrm{O}(d)$  $d$ 
\begin{equation}
  \so(d) = \{ \Omega \in \R^{d \times d} : \Omega^\top = -\Omega \},
\end{equation}
 $d \times d$ $\so(d)$ 
$\frac{d(d-1)}{2}$ \textbf{}

\textbf{Cartan } $g \in \mathrm{O}(d)$ 

\begin{equation}
  g = \exp(\xi), \quad \xi \in \so(d),
\end{equation}
 $\exp: \so(d) \to \mathrm{O}(d)$ 
$\log: \mathrm{O}(d) \to \so(d)$ $\mathrm{O}(d)$ 
 $\det(g) = -1$  $g$  $-1$ 
\end{definition}

\bigskip
\noindent
\textbf{Definition (Cartan Decomposition).}
Let $G = \mathrm{O}(d)$ be the $d$-dimensional orthogonal group with Lie algebra:
\begin{equation}
  \so(d) = \{ \Omega \in \R^{d \times d} : \Omega^\top = -\Omega \},
\end{equation}
the set of all $d \times d$ skew-symmetric matrices. $\so(d)$ is a
\textbf{vector space} of dimension $\frac{d(d-1)}{2}$.

\textbf{Cartan decomposition} states: every $g \in \mathrm{O}(d)$ near the
identity can be uniquely written as:
\begin{equation}
  g = \exp(\xi), \quad \xi \in \so(d),
\end{equation}
where $\exp: \so(d) \to \mathrm{O}(d)$ is the matrix exponential. Its
inverse $\log: \mathrm{O}(d) \to \so(d)$ (the matrix logarithm) is
defined everywhere on $\mathrm{O}(d)$ except a measure-zero subset
where $\det(g) = -1$ and $g$ has eigenvalue $-1$.

% -------------------------------------------------------------------
\subsection{}
\subsection*{Linearization of the Gauge-Fixing Problem}

\noindent
\textbf{.}
 ACE  $\{g_m\}_{m=1}^M \subset \mathrm{O}(d)$
 $M$ 
$\mathcal{L}(\{g_m\})$

\textbf{}

\begin{enumerate}[label=\textbf{ \arabic*}:, leftmargin=*]
  \item  $g_0 \in \mathrm{O}(d)$
  \item  $g_m$ 
        \begin{equation}\label{eq:param}
          g_m = g_0 \cdot \exp(\xi_m), \quad \xi_m \in \so(d).
        \end{equation}
  \item  $\xi_m$ $\so(d)$  $\R^{d(d-1)/2}$
         $\{E_{ij} = e_i e_j^\top - e_j e_i^\top\}_{1 \leq i < j \leq d}$
        \begin{equation}\label{eq:basis}
          \xi_m = \sum_{1 \leq i < j \leq d} \omega_m^{ij} \cdot E_{ij},
          \quad \omega_m^{ij} \in \R.
        \end{equation}
  \item  $\mathcal{L}$  $\{\omega_m^{ij}\}$ 
        \begin{equation}\label{eq:expand}
          \mathcal{L}(\{\omega_m^{ij}\}) \approx
          \mathcal{L}_0 + \mathbf{J} \cdot \bm{\omega} + \frac{1}{2} \bm{\omega}^\top \mathbf{H} \bm{\omega},
        \end{equation}
         $\bm{\omega} \in \R^{M \cdot d(d-1)/2}$ 
        $\mathbf{J}$ Jacobian$\mathbf{H}$  Hessian
  \item  $\mathbf{H} \cdot \bm{\omega}^* = -\mathbf{J}^\top$ 
        Newton  $\exp$ 
\end{enumerate}

\textbf{}  2--5 \textbf{}
$\R^{M \cdot d(d-1)/2}$  $\mathrm{O}(d)^M$ 
 Cartan ——


\medskip
\noindent
\textbf{English.}
Consider the gauge-fixing problem in ACE potentials. Let
$\{g_m\}_{m=1}^M \subset \mathrm{O}(d)$ be the gauge transformations
(rotation matrices) of $M$ experts. The gauge-fixing objective is to
minimize a loss function $\mathcal{L}(\{g_m\})$ measuring inconsistency
among rotated potentials.

\textbf{Linearization procedure}:

\begin{enumerate}[label=\textbf{Step \arabic*}:, leftmargin=*]
  \item Choose a reference rotation $g_0 \in \mathrm{O}(d)$ (e.g., the current estimate).
  \item Parameterize each $g_m$ as:
        \begin{equation}
          g_m = g_0 \cdot \exp(\xi_m), \quad \xi_m \in \so(d).
        \end{equation}
  \item Vectorize the skew-symmetric matrix $\xi_m$. $\so(d)$ is isomorphic to
        $\R^{d(d-1)/2}$ with basis
        $\{E_{ij} = e_i e_j^\top - e_j e_i^\top\}_{1 \leq i < j \leq d}$:
        \begin{equation}
          \xi_m = \sum_{1 \leq i < j \leq d} \omega_m^{ij} \cdot E_{ij},
          \quad \omega_m^{ij} \in \R.
        \end{equation}
  \item Expand the loss $\mathcal{L}$ to first order in $\{\omega_m^{ij}\}$:
        \begin{equation}
          \mathcal{L}(\{\omega_m^{ij}\}) \approx
          \mathcal{L}_0 + \mathbf{J} \cdot \bm{\omega} + \frac{1}{2} \bm{\omega}^\top \mathbf{H} \bm{\omega},
        \end{equation}
        where $\bm{\omega} \in \R^{M \cdot d(d-1)/2}$ concatenates all parameters,
        $\mathbf{J}$ is the gradient (Jacobian), and $\mathbf{H}$ is the approximate Hessian.
  \item Solve the linear system $\mathbf{H} \cdot \bm{\omega}^* = -\mathbf{J}^\top$
        for the Newton step, then map back to the group manifold via $\exp$.
\end{enumerate}

\textbf{Key insight:} All computations in Steps 2--5 take place in the
\textbf{vector space} $\R^{M \cdot d(d-1)/2}$, not on the curved group
manifold $\mathrm{O}(d)^M$. This is the power of Cartan decomposition —
transforming an optimization problem on a non-abelian group into a standard
linear least-squares problem on a vector space.

% -------------------------------------------------------------------
\subsection{}
\subsection*{Complexity Analysis}

\begin{proposition}[]
\label{prop:complexity}

 $M$  $d$ 
\begin{itemize}
  \item $M \cdot \frac{d(d-1)}{2}$
  \item  $\exp(\xi_m)$ $\mathcal{O}(d^3)$
        Rodrigues  $\mathcal{O}(d^2)$ $d=3$  SO(3)
  \item Jacobian $\mathcal{O}(M \cdot d^3)$
  \item $\mathcal{O}((M d^2)^3)$ 
         $\mathcal{O}(M d^2 \cdot \kappa)$ 
\end{itemize}

 ($d=3$, $M \lesssim 100$) $\mathcal{O}(M \cdot 27)$
—— Riemannian 
\end{proposition}

\bigskip
\noindent
\textbf{Proposition (Computational Overhead of Linearization).}
For $M$ experts and $d$-dimensional rotation group:
\begin{itemize}
  \item Total parameters: $M \cdot \frac{d(d-1)}{2}$.
  \item Matrix exponential $\exp(\xi_m)$ per iteration: $\mathcal{O}(d^3)$
        (reducible to $\mathcal{O}(d^2)$ via Rodrigues formula, especially
        for $d=3$ i.e., SO(3)).
  \item Jacobian computation: $\mathcal{O}(M \cdot d^3)$ via automatic
        differentiation or analytic gradients.
  \item Linear system solve: $\mathcal{O}((M d^2)^3)$ for dense direct
        solver, or $\mathcal{O}(M d^2 \cdot \kappa)$ for conjugate gradient.
\end{itemize}

For practical settings ($d=3$, $M \lesssim 100$), total overhead is
$\mathcal{O}(M \cdot 27)$ floating-point operations — entirely feasible
and far cheaper than full Riemannian optimization.

% ===================================================================
\section{Riemannian Gauss-Newton }
\section*{When Riemannian Gauss-Newton IS Needed}

\noindent
\textbf{.}
Cartan  $\log: \mathrm{O}(d) \to \so(d)$ 


\begin{enumerate}[label=(\arabic*)]
  \item \textbf{:}  $g_m$  $> \pi/2$
        $\log(g_m)$ 
  \item \textbf{:} $\mathrm{O}(d)$  $\det = -1$ 
        
  \item \textbf{:}  $\mathrm{O}(d)$ 
        Riemannian  Riemannian trust-region 
\end{enumerate}

\textbf{Riemannian Gauss-Newton }——

 Riemannian 

\medskip
\noindent
\textbf{English.}
Cartan linearization is not a panacea. The logarithm map
$\log: \mathrm{O}(d) \to \so(d)$ is valid only under:

\begin{enumerate}[label=(\arabic*)]
  \item \textbf{Proximity to identity:} When $g_m$ is far from the identity
        (e.g., rotation angle $> \pi/2$), $\log(g_m)$ may not exist or may
        incur large linearization error;
  \item \textbf{Reflection components:} $\mathrm{O}(d)$ includes the
        $\det = -1$ reflection branch, where the logarithm has singularities;
  \item \textbf{Large-scale global optimization:} When searching for a global
        optimum over all of $\mathrm{O}(d)$, Riemannian optimization (e.g.,
        Riemannian trust-region methods) is more robust.
\end{enumerate}

\textbf{In these cases, Riemannian Gauss-Newton is the appropriate tool} —
but its convergence remains linear (when residuals are non-zero),
and its per-iteration cost far exceeds that of the linearized method
(requiring computation of Riemannian connections, parallel transport,
and geodesics).

% ===================================================================
\section{ \\ Before-and-After Comparison}
% ===================================================================

\begin{table}[h]
\centering
\caption{5 \\ Comparison Before and After Correction}
\begin{tabular}{@{}p{4.5cm}p{5cm}p{5cm}@{}}
\toprule
\textbf{\newline Aspect} &
\textbf{\newline Original} &
\textbf{\newline Corrected} \tabularnewline
\midrule

\newline Gauge-fixing method &
``''\newline ``Non-abelian gauge fixing'' &
Cartan  $\to$ \newline Cartan decomp. $\to$ vector-space linear LS \tabularnewline
\midrule

\newline Mathematical mechanism &
Yang-Mills Gribov\newline Yang-Mills style (ghosts, Gribov) &
 $\log: \mathrm{O}(d) \to \so(d)$\newline Logarithm map $\log: \mathrm{O}(d) \to \so(d)$ \tabularnewline
\midrule

\newline Computational complexity &
\newline Unspecified, implied high &
$\mathcal{O}(d^2)$ \newline $\mathcal{O}(d^2)$ overhead, near-zero near identity \tabularnewline
\midrule

Gauss-Newton \newline Gauss-Newton convergence &
\newline Claimed quadratic &
\textbf{}\newline \textbf{Linear convergence} (non-zero residuals) \tabularnewline
\midrule

\newline Scope of applicability &
\newline Unqualified &
 Riemannian\newline Linearized near identity; Riemannian when far \tabularnewline
\bottomrule
\end{tabular}
\end{table}

% ===================================================================
\section{ \\ Summary}
% ===================================================================

\noindent
\textbf{:}

5——ACE  O(3)  O(d)——\textbf{}
``''\textbf{}

\begin{enumerate}[label=(\arabic*)]
  \item \textbf{Cartan :}  $\log: \mathrm{O}(d) \to \so(d)$
         $\so(d) \cong \R^{d(d-1)/2}$
        ——Gribov  BRST 
  \item \textbf{:} ``''
        Faddeev-Popov  ACE 
        
  \item \textbf{:} Gauss-Newton 
        ACE 
\end{enumerate}

 Cartan 
 $\mathcal{O}(d^2)$
Riemannian 

\bigskip
\noindent
\textbf{English Summary:}

The core physical intuition of Viewpoint 5 — that ACE potentials have O(3)
rotational symmetry, making the gauge group O(d) — is \textbf{correct}.
However, concluding that this requires ``non-abelian gauge fixing'' is
\textbf{overstated} for three reasons:

\begin{enumerate}[label=(\arabic*)]
  \item \textbf{Existence of Cartan decomposition:} The logarithm map
        $\log: \mathrm{O}(d) \to \so(d)$ linearizes the problem on the
        group manifold to the vector space $\so(d) \cong \R^{d(d-1)/2}$,
        reducing it to standard linear least squares — no ghost fields,
        Gribov copies, or BRST symmetry required.
  \item \textbf{Misuse of physics terminology:} ``Non-abelian gauge fixing''
        has a precise technical meaning in quantum field theory
        (Faddeev-Popov procedure), and none of that machinery applies to
        the finite-dimensional parameter estimation problem of ACE potentials.
  \item \textbf{Incorrect convergence claim:} Gauss-Newton has quadratic
        convergence only in the zero-residual case. The ACE gauge-fixing
        problem has non-zero residuals, so convergence is only linear.
\end{enumerate}

The corrected formulation accurately reflects the true mathematical structure:
a non-abelian symmetry that can be linearized via Cartan decomposition, with
$\mathcal{O}(d^2)$ solution complexity, requiring full Riemannian optimization
only when far from the solution.

% ===================================================================
\section*{ \\ Acknowledgments}
% ===================================================================

\noindent
 O(d)  Gauss-Newton



\bigskip
\noindent
\textit{We thank the reviewer for the penetrating insight into the relationship
between the O(d) gauge group and its Lie algebra, and for the accurate reminder
about Gauss-Newton convergence theory. This correction is based on standard
results in numerical optimization, Lie group theory, and gauge fixing in
quantum field theory.}

\end{document}
